\chapter{Corps algébriquement clos}

Dans cette partie, toutes les structures sont des anneaux considérés comme
$\cL_{ring}$-structures donc plutôt que de noter $\cA_{ring}$ ou $\cA$, on
désignera par la même lettre $A$ à la fois le modèle et le domaine. 

Pour des raisons de convenances, on rappelle ici le langage et les axiomes pour
les anneaux et les corps qui avaient été donné plus haut.

\begin{defi}[Anneaux, corps]
Le language naturel des anneaux est $\cL_{ring} = (+,\cdot,-,0,1)$.

On a les axiomes des groupes abéliens (dans le language des groupes additifs
$\cL_{agp} =(+,-,0))$

\begin{enumerate}[label=\textbf{R\arabic*}]
    \item $\forall x,y,z [(x+y)+z = x+(y+z)]$,
    \item $\forall x [x+0 \wedge 0+x = x]$,
    \item $\forall x [x+(-x)=0 \wedge (-x)+x=0]$,
    \item $\forall x,y [x+y = y+x]$.
\end{enumerate}

Ensuite les axiomes de monoïde commutatif (dans le language $\cL_{mon} = (\cdot,
1))$

\begin{enumerate}[label=\textbf{R\arabic*},start=5]
    \item $\forall x,y,z [(x\cdot y)\cdot z = x\cdot(y\cdot z)]$,
    \item $\forall x [x\cdot 1 \wedge 1\cdot x = x]$,
    \item $\forall x,y [x\cdot y = y\cdot x]$.
\end{enumerate}

et la distributivité

\begin{enumerate}[label=\textbf{R\arabic*},start=8]
    \item $\forall x,y,z [x \cdot (y+z) = (x \cdot y) + (x \cdot z)]$.
\end{enumerate}

Ces huit phrases axiomatisent la théorie des anneaux commutatifs.

On rajoute ensuite les axiomes 

\begin{enumerate}[label=\textbf{F\arabic*}]
    \item $0 \neq 1$,
    \item $\forall x [x = 0 \vee \exists y [x\cdot y = 1]]$.
\end{enumerate}

pour obtenir les axiomes de la théorie des corps.

\end{defi}

\section{Propriétés de base}

On rappelle ensuite quelques définitions et propriétés qui nous seront utiles
par la suite. On donne certaines démonstrations, on pourra trouver les autres
dans n'importe quel livre d'algèbre, par exemple dans \cite{Berhuy_2025}.

\begin{defi}[Anneau intègre]
Un anneau est \kw{intègre} si il est sans diviseur de zéro, c'est à dire si
c'est un modèle pour la phrase \[ \forall x,y [x \cdot y = 0 \rightarrow (x=0
\vee y=0) ] \].
\end{defi}

\begin{prop}
    Si $A$ est un anneau intègre, il possède un corps des fractions $K_A$ qui
    vérifie la propriété universelle suivante:
    Pour tout plongement $f : A \inj{} F$ de $A$ dans un corps $F$, il existe un
    unique plongement $\tilde{f} : K_A \inj{} F$ tel que le diagramme suivant commute.

    % https://tikzcd.yichuanshen.de/#N4Igdg9gJgpgziAXAbVABwnAlgFyxMJZABgBpiBdUkANwEMAbAVxiRAEEQBfU9TXfIRQBGclVqMWbANIB9Tjz7Y8BIqOHj6zVohAAxbuJhQA5vCKgAZgCcIAWyRkQOCElEgGdAEYwGABX4VIQ8YSxwQai0pXSwI5zosBjYACwgIAGtuXhAbezdqFyQAJmoGLDAdECg6OGTjOM8ff0DBNgZQ8MjJSoAdHpgADyw4HDgAAgBCPrwGWGBLLjicBKTdVIysq1sHRCdCxBKJbTZLBu9fAOVW3WssE2TO+MSUtMyuCi4gA
    \[
    \begin{tikzcd}
    A \arrow[r, "i", hook] \arrow[rd, "f"', hook] & K_A \arrow[d, "\exists !\tilde{f}", dashed, hook] \\
                                                & F                                                
    \end{tikzcd}
    \]
\end{prop}

\begin{exem}
On note $A[X_1, \ldots, X_n]$ l'anneau des polynômes à coefficients dans $A$ en
$n$ indéterminées dans un anneau $A$. Si $A$ est intègre, $A[X_1, \ldots, X_n]$
également. Si $A$ est un corps, le corps des fractions de $A[X_1, \ldots, X_n]$
est noté $A(X_1, \ldots, X_n)$ et on l'appelle le corps des fractions
rationnelles.
\end{exem}

\todo[inline]{defs: Morphisme d'anneau,noyau, idéal, idéal engendré,idéal
propre, idéal premier, idéal principal, idéal maximal, caractéristique,
sous-corps premier
}

\todo[inline]{props: noyau d'un morphisme est un idéal, quotient par un idéal, 
R/I intègre ssi I premier, R/I corps ssi I maximal,}

\section{Extensions de corps}

\todo[inline]{defs : extension, extension simple, extension finie, anneau
principal, polynome minimal, extension algébriquement,
 }
\todo[inline]{prop: th de la base de Hilbert,si F c K c L, K algébrique sur F et L algébrique sur K alors L algébrique sur F, nombre de racines dans K[X]}

\section{Cloture algébrique d'un corps}

\section{Catégoricité, complétude}

\section{Ensembles définissables, variétés}